%% kapa: exam=2 round=28 year=2017 period=2 subject=appraisal_theory subject_ko=감정평가이론 seq=1 number=1 points=40
%% source: Q-Net 감정평가사 기출문제 (cst00309) artlSeq=5201712 file=raw/1008/5201712/2017년 감정평가사 2차시험 기출문제(감정평가 이론).pdf
%% conversion: kordoc 4.24.51 -> tools/gen2.py (auto); review=proofread-sonnet; flags=none
\begin{longquestion}{【문제 1】}
제시된 자료를 참고하여 다음 물음에 답하시오. (40점) \par
\begin{framedtext}
감정평가사 甲은 감정평가사 乙이 작성한 일반상업지역 내 업무용 부동산(대지면적: 3,000m\textsuperscript{2}, 건물: 30년 경과된 철근콘크리트조 6층)에 대한 감정평가서를 심사하고 있다. 동 감정평가서에 따르면, 인근지역은 일반적으로 대지면적 200m\textsuperscript{2} \textasciitilde{} 500m\textsuperscript{2} 내외 2층 규모의 상업용으로 이용되고 있으며, 최근 본건 부동산 인근에 본건과 대지면적이 유사한 토지에 20층 규모의 주거 및 상업 복합용도 부동산이 신축되어 입주(점) 중에 있는 것으로 조사되어 있다. 검토결과 원가방식(면적 400m\textsuperscript{2} 상업용 나대지의 최근 매매사례 단가를 적용한 토지가치에 물리적 감가수정만을 행한 건물가치 합산)에 의한 시산가치가 수익방식(현재 본건 계약임대료 기준)에 의한 시산가치보다 높게 산출되어 있다.
\end{framedtext}
\medskip\noindent 1) 심사 감정평가사 甲은 감정평가사 乙에게 추가적으로 최유효이용 분석을 요청하였는 바, 최유효이용 판단 기준을 설명하고, 구체적인 최유효이용 분석방법을 설명하시오. (20점) \par
\medskip\noindent 2) 최유효이용에 대한 두 가지 분석 유형(방법)에 따른 결과가 다르다면, 그 이유와 그것이 의미하는 바를 설명하시오. (10점) \par
\medskip\noindent 3) 원가방식에 의한 시산가치가 수익방식에 의한 시산가치보다 높게 산출된 것이 타당한 것인지 감정평가 원리(원칙)를 기준으로 설명하고, 올바른 원가방식 적용방법에 관하여 설명하시오. (10점) \par
\end{longquestion}
